%--------------------------------------------------------------------
\chapter*{Prerequisiti matematici}
\addcontentsline{toc}{chapter}{Prerequisiti matematici}
\setlength{\parindent}{2pt}

\begin{multicols*}{2}

    \section*{Analisi matematica}
    \addcontentsline{toc}{section}{Analisi matematica}

    \begin{itemize}
        \item \textbf{Teorema della funzione implicita} -- Sia $U$ un aperto di $\RR^m \times \RR^n$ e sia
              $f : U \to \RR^n$ una funzione di classe $C^k$ con $k \geq 1$. Sia $\vec{p} = (\vec{x_0}, \vec{y_0})$ un punto in $U$ con
              $f(\vec{x_0}, \vec{y_0}) = \vec{a}$ e $J_{\vec{y}} f(\vec{p})$ invertibile. \smallskip

              Allora esiste un
              intorno $A = I_{\vec{x}} \times I_{\vec{y}}$  di $\vec{p}$ in $U$ all'interno del quale esiste un'unica
              funzione $g : I_{\vec{x}} \to I_{\vec{y}}$ di classe $C^k$ per cui:
              \[ \vec{y} = g(\vec{x}) \iff f(\vec{x}, \vec{y}) = \vec{a}, \quad \text{(in $A$)}. \]
              Inoltre per tale $g$ vale:
              \[ J g(\vec{x_0}) = - J_{\vec{y}} f(\vec{p})\inv J_{\vec{x}} f(\vec{p}). \]

        \item \textbf{Teorema di invertibilità locale} -- Sia $U$ un aperto di $\RR^n$ e sia $f : U \to \RR^n$ una
              funzione di classe $C^k$, con $k \geq 1$. Sia $\vec{x_0}$ un punto in $U$ con $J f(\vec{x_0})$ invertibile. \smallskip

              Allora esiste un
              intorno $A$ di $\vec{x_0}$ in $U$ dentro al quale $\restr{f}{A}$ ha un'inversa $g$, anch'essa di classe
              $C^k$, per la quale $J g(f(\vec{x})) = J f(\vec{x})\inv$.

        \item \textbf{Teorema di esistenza e unicità globale per sistemi lineari di equazioni differenziali} --
              Sia $I \subseteq \RR$ un intervallo aperto e siano date due funzioni continue:
              \[ A : I \to \RR^{n \times n} \quad \text{e} \quad \vec{b} : I \to \RR^n. \]
              Fissato $(t_0, \vec{y_0}) \in I \times \RR^n$, esiste un'unica soluzione $\vec{y} : I \to \RR^n$
              del problema di Cauchy:
              \[
                  \begin{cases}
                      \vec{y}'(t) = A(t)\vec{y}(t) + \vec{b}(t), \\
                      \vec{y}(t_0) = \vec{y_0}.
                  \end{cases}
              \]

        \item \textbf{Teorema di Cauchy-Lipschitz per l'esistenza e l'unicità locale} --
              Sia $\Omega$ un aperto di $\RR \times \RR^n$ e sia $\vec{f} : \Omega \to \RR^n$ una funzione continua.
              Si supponga inoltre che $\vec{f}$ sia \emph{localmente lipschitziana} rispetto alla seconda variabile. \smallskip

              Allora, per ogni $(t_0, \vec{y_0}) \in \Omega$, esistono $\delta > 0$ e un'unica funzione
              $\vec{y} : (t_0 - \delta, t_0 + \delta) \to \RR^n$ di classe $C^1$ che risolve il problema di Cauchy:
              \[
                  \begin{cases}
                      \vec{y}'(t) = \vec{f}(t, \vec{y}(t)), \\
                      \vec{y}(t_0) = \vec{y_0}.
                  \end{cases}
              \]

        \item \textbf{Teorema di dipendenza liscia dai dati iniziali} --
              Sia $\Omega$ un aperto di $\RR \times \RR^n$ e sia $\vec{f} : \Omega \to \RR^n$ una funzione di classe
              $C^k$ con $k \geq 1$. Indichiamo con $\Phi(t, t_0, \vec{y_0})$ la soluzione massimale del
              problema di Cauchy con dato iniziale $\vec{y}(t_0) = \vec{y_0}$. \smallskip

              Allora l'insieme di definizione del flusso:
              \[
                  \mathcal{D} = \left\{ (t, t_0, \vec{y_0}) \in \RR \times \Omega :
                  \begin{array}{c}
                      \text{la soluzione } \vec{y}(t, t_0, \vec{y_0}) \\
                      \text{esiste al tempo } t
                  \end{array}
                  \right\}
              \]
              è un aperto di $\RR \times \Omega$ e l'applicazione $\Phi : \mathcal{D} \to \RR^n$ è di classe $C^k$.
    \end{itemize}

\end{multicols*}